\subsection{半单李代数的定义}

定义 1。设 \(\mathfrak{g}\) 是一个李代数。当且仅当 \(\mathfrak{g}\) 的唯一交换理想是 \(\mathfrak{g}\) 时，称 \(\{0\}\) 为半单的。

注记。（1）零代数 \(\{0\}\) 是半单的。维数为 1 或 2 的代数不是半单的（参见~§ 5，第 1 号，例 1）。存在维数为 3 的半单代数（参见第 7 号）。

\begin{enumerate}
\def\labelenumi{(\arabic{enumi})}
\setcounter{enumi}{1}
\item
  半单代数的中心为零，因而其伴随表示是忠实的。
\item
  若 \(g_{1}, \ldots, g_{n}\) 是半单的，则 \(g = g_{1} \times \cdots \times g_{n}\) 是半单的：因为若 a 是 g 的交换理想，a 在 \(g_{1}, \ldots, g_{n}\) 上的投影都退化为 \(\{0\}\)。
\end{enumerate}

定理 1。设 \(\mathfrak{g}\) 是一个李代数。以下条件等价：

\begin{enumerate}
\def\labelenumi{(\alph{enumi})}
\item
  g 是半单的。
\item
  g 的根基 r 为零。
\item
  \(\beta\) 的 Killing 形式 \(\mathfrak{g}\) 非退化。
\end{enumerate}

此外，半单李代数等于其导出理想。

\begin{enumerate}
\def\labelenumi{(\alph{enumi})}
\item
  \(\Rightarrow\) (b)：因为若 \(\mathfrak{r} \neq \{0\}\)，则 \(\mathfrak{r}\) 的最后一个非零导出代数是 \(\mathfrak{g}\) 的交换理想。
\item
  \(\Rightarrow\) (c)：这由 § 5，第 5 号的命题 5 (b) 得到（这同时证明了定理的最后一个断言）。
\item
  \(\Rightarrow\) (a)：这由 §4，第 4 号的命题 6 (b) 得到。
\end{enumerate}

推论。设 \(\mathfrak{g}\) 是半单李代数，\(\rho\) 是 \(\mathfrak{g}\) 在有限维向量空间 V 上的表示。则 \(\rho(\mathfrak{g}) \subset \mathfrak{sl}(V)\)。

线性形式 \(x \mapsto \operatorname{Tr} \rho(x)\)（\(x \in \mathfrak{g}\)）当 \(x\) 形如 \([y, z]\)（\(y \in \mathfrak{g}, z \in \mathfrak{g}\)）时为零，因而在 \(\mathcal{D}\mathfrak{g} = \mathfrak{g}\) 上也为零。

命题 1。设 \(\mathfrak{g}\) 是半单李代数，\(\rho\) 是 \(\mathfrak{g}\) 的有限维忠实表示。则 \(\mathfrak{g}\) 所关联的 \(\rho\) 上的双线性形式非退化。

该形式关于 g 的正交子空间是一个可解理想（§ 5，第 4 号，定理 2），因而为零。

推论 1。设 g 是一个李代数，\(\beta\) 是其 Killing 形式，\(\mathfrak{a}\) 是 g 的半单子代数。关于 \(\mathfrak{b}\) 与 \(\mathfrak{a}\) 正交的子空间 \(\beta\) 是 g 中 \(\mathfrak{a}\) 的补空间，且 \({[}\mathfrak{a}, \mathfrak{b}{]} \subset \mathfrak{b}\)。若 \(\mathfrak{a}\) 是 g 的理想，则 \(\mathfrak{b}\) 也是理想，此时它是 g 中 \(\mathfrak{a}\) 的中心化子。

设 \(\beta'\) 是 \(\beta\) 在 \(a\) 上的限制：它是 \(x \mapsto \operatorname{ad}_{\mathfrak{g}} x\) 在空间 \(a\) 上由表示 \(\mathfrak{g}\) 所关联的双线性形式。该表示是忠实的，因而 \(\beta'\) 非退化（命题 1）。所以 \(\mathfrak{b}\) 是 \(a\) 在 \(\mathfrak{g}\) 中的补空间。另一方面，若 \(x, y\) 属于 \(a\) 且 \(z \in \mathfrak{b}\)，则

\[
\beta (x, [ y, z ]) = \beta ([ x, y ], z) = 0,
\]

由于 \([x,y] \in \mathfrak{a}\)，所以 \([y,z] \in \mathfrak{b}\)，这证明了 \([\mathfrak{a},\mathfrak{b}] \subset \mathfrak{b}\)。若 \(\mathfrak{a}\) 是 \(\mathfrak{g}\) 的理想，我们知道 \(\mathfrak{b}\) 也是 \(\mathfrak{g}\) 的理想（§ 3，命题 7），且 \(\mathfrak{g}\) 可与 \(\mathfrak{a} \times \mathfrak{b}\) 等同。由于 \(\mathfrak{a}\) 的中心为零，\(\mathfrak{a}\) 中 \(\mathfrak{g}\) 的中心化子是 \(\mathfrak{b}\)。

推论 2。半单李代数的每个半单李代数扩张都是半单的且平凡。

这立即由推论 1 得到。

推论 3。若 g 是半单的，则 g 的每个导子都是内导子。

ad g 同构于 g，因而是半单的，并且是李代数 d（由 g 的导子组成）的理想（§ 1，命题 1）。若 D ∈ d 与 ad g 的元素可交换，则对所有 x ∈ g，有 ad D(x) = {[}D, ad x{]} = 0，从而 D(x) = 0；所以 D = 0。推论 3 现在由推论 1 得到。
% semantic-fragments: legacy-input-seam
\relax{}
